\chapter{Trabajo pendiente y futuro}

El trabajo inmediato consiste en completar el ensamble y la caracterización
física, medir geometría y masas, identificar los servos sustituidos y verificar
holguras y topes. Después se debe cerrar el diseño eléctrico, implementar una
parada física mediante OE, calibrar cada MG996R y validar progresivamente un
servo, una pata y el robot suspendido.

La marcha nominal física se transferirá solo cuando se cumplan esas condiciones.
Las pruebas avanzarán desde postura con soporte hasta un paso, un ciclo de
gateo y marcha continua, registrando orientación, potencia, intervenciones y
desplazamiento.

En simulación se investigará una liberación trasera que produzca una pérdida
filtrada sostenida sin aumentar la altura temprana por encima de 0,80 ni superar
un salto articular de 0,20~rad. El supervisor deberá someterse a pruebas
provocadas de pérdida de datos, contacto inconsistente y margen negativo antes
de habilitar estas paradas.

Finalmente se formulará y entrenará la política PPO como una capa de
correcciones acotadas. La comparación final conservará las semillas y tamaños
de muestra definidos y evaluará la misma marcha con y sin corrección aprendida.
